We consider a discrete-time timetabling problem on a weekly grid.
The entity/constraint structure is aligned with curriculum-based course timetabling conventions used in ITC-style benchmarks~\cite{bonutti2012benchmarking,bettinelli2015overview}.

\paragraph{Calendar.}
Let $\mathcal{W}$ be the set of teaching weeks, $\mathcal{D}$ the set of teaching days, and $\mathcal{S}=\{0,\dots,S-1\}$ the set of within-day time slots.
We assume a fixed set of teaching days (Monday--Saturday) and exclude Sunday.

\paragraph{Entities.}
We are given:
\begin{itemize}[leftmargin=*]
  \item groups $g \in \mathcal{G}$ (student cohorts), each with a size $\mathrm{size}(g)$,
  \item staff $p \in \mathcal{P}$ (professors and TAs), each with
    day-availability $\mathrm{AvailDays}(p)\subseteq \mathcal{D}$ and optional
    week-availability $\mathrm{AvailWeeks}(p)\subseteq\mathcal{W}$,
  \item rooms $r \in \mathcal{R}$ with type $\mathrm{type}(r)$ and capacity $\mathrm{cap}(r)$,
    where capacity can be configured either categorically (\texttt{SMALL/MEDIUM/BIG})
    or numerically (exact seat counts),
  \item activities $a \in \mathcal{A}$, each with:
    \begin{itemize}[leftmargin=*]
      \item week $\mathrm{week}(a)\in\mathcal{W}$,
      \item kind $\mathrm{kind}(a)\in\{\mathrm{LEC},\mathrm{TUT},\mathrm{LAB}\}$,
      \item duration $\mathrm{dur}(a)\in\{1,2,3\}$ (number of consecutive slots),
      \item participating groups $\mathcal{G}(a)\subseteq \mathcal{G}$,
      \item assigned staff $\mathrm{staff}(a)\in\mathcal{P}$,
      \item optional specialization tag requirement $\mathrm{tag}(a)$ for specialized labs.
    \end{itemize}
\end{itemize}

\paragraph{Decision.}
For each activity $a$, choose a start day and start slot $(d,s)$ in week $\mathrm{week}(a)$; optionally choose a room $r$.

\paragraph{Feasibility.}
Hard constraints include resource non-overlap (groups, staff, rooms), staff availability, room eligibility, and several rule-based constraints such as ``week 1 is lectures-only'' and optional staff load caps.

\paragraph{Quality.}
We also consider soft constraints expressed as a weighted penalty:
free-day shortfalls, gaps, late starts, stability, room consistency,
and an additional weekly concentration term discouraging many repeated
same-course LEC/TUT sessions for one group in one week (useful during manual week-shift editing).
